%============================================================================
\section{Fundamento teórico}\label{sec:fundamento}
%============================================================================

\subsection{Modulación BPSK}\label{subsec:bpsk}

En BPSK cada bit se asocia a uno de dos símbolos antipodales. En banda base, la
señal con pulso de conformación $p(t)$ es
\begin{equation}
  s(t)=A\sum_{k} a_k\,p(t-kT_s), \qquad a_k\in\{-1,+1\},
  \label{eq:bpsk}
\end{equation}
donde $A$ es la amplitud y $T_s=1/R_s$ el periodo de símbolo. La información
está solo en la componente en fase (I); con $M=2$, $k=\log_2M=1$ y $R_b=R_s$.
Con pulso rectangular, $E_b=A^2T_s$, los símbolos están separados $2\sqrt{E_b}$
y, en AWGN con filtro acoplado y umbral en cero~\cite{proakis2008},
\begin{equation}
  P_b^{\mathrm{AWGN}}=Q\!\left(\sqrt{2E_b/N_0}\right).
  \label{eq:pb_awgn}
\end{equation}
Su PSD es $S_s(f)=A^2T_s\,\mathrm{sinc}^2(fT_s)$, con nulos en $f=\pm mR_s$, de
modo que el ancho de banda entre primeros nulos es $2R_s$~\cite{sklar2001}. El
lóbulo principal contiene cerca del \SI{90}{\percent} de la potencia y el primer
lóbulo lateral está unos \SI{13,3}{\decibel} por debajo; con un pulso RRC de
factor $\beta$, el ancho de banda sería $(1+\beta)R_s$.

\textbf{Alternativa descartada.} En OOK, con igual energía por bit, los símbolos
$0$ y $\sqrt{2E_b}$ están un factor $\sqrt2$ más cerca
($P_b=Q(\sqrt{E_b/N_0})$, \SI{3}{\decibel} de penalización)~\cite{proakis2008}, y
su umbral depende del nivel recibido, por lo que un impulso lo cruza con más
facilidad. Por ello se eligió BPSK.

\subsection{Ruido impulsivo frente a AWGN}\label{subsec:ruido}

El AWGN es gaussiano, estacionario y afecta por igual a todas las muestras. La
interferencia de origen humano es no gaussiana: los modelos de Middleton la
describen como una mezcla dominada por pocos eventos intensos~\cite{middleton1977},
y en tiempo discreto se usa el modelo Bernoulli-gaussiano~\cite{ghosh1996}. En
este trabajo se empleó un tren periódico controlado de impulsos rectangulares
de amplitud $A_i$, ancho $W$ y periodo $N_T$ muestras,
\begin{equation}
  i[n]=A_i\sum_{k}\Pi_W[n-kN_T],
  \label{eq:tren}
\end{equation}
donde $\Pi_W[n]=1$ para $0\le n<W$. Representa las componentes periódicas del
ruido de conmutación~\cite{zimmermann2002} y permite variar solo la amplitud.
Su ciclo útil es $\delta=W/N_T$ y su potencia media $A_i^2\delta$: con $\delta$
pequeño, la energía se concentra en pocas muestras y se reparte en una banda
ancha, con envolvente espectral de primer nulo en $f_s/W$ y líneas espaciadas
$f_T=f_s/N_T$; al ser unipolar, el tren también tiene una componente de continua
$A_i\delta$. La Tabla~\ref{tab:awgn_vs_imp} resume las diferencias con el AWGN.

\begin{table}[!t]
\centering
\caption{Comparación entre AWGN y ruido impulsivo periódico}
\label{tab:awgn_vs_imp}
\footnotesize
\setlength{\tabcolsep}{3pt}
\begin{tabular}{@{}p{0.24\columnwidth}p{0.35\columnwidth}p{0.35\columnwidth}@{}}
\toprule
\textbf{Aspecto} & \textbf{AWGN} & \textbf{Impulsivo periódico} \\
\midrule
Distribución temporal & Continua, todas las muestras & Concentrada en $W$ de cada $N_T$ muestras \\
Amplitud & Gaussiana, media cero & Pico $A_i$, alta relación pico/media \\
PSD & Piso plano & Líneas cada $f_T$ con envolvente ancha \\
Constelación & Nube circular en ambos símbolos & Puntos aislados desplazados $A_i$ \\
Error en BPSK & Gradual con $E_b/N_0$ & Umbral en $A_i W/\mathrm{SPS}\approx A$; piso de error \\
\bottomrule
\end{tabular}
\end{table}

\begin{figure}[!t]
\centering
\captura[0.95\columnwidth]{escenario.png}
\caption{Escenario S1: nodo Tx sobre la máquina, estación Rx de supervisión y
fuentes de perturbación impulsiva (motores, contactores y variadores).
Ilustración elaborada con apoyo de IA (Claude, Anthropic), que generó el código
de dibujo; el contenido técnico fue revisado por el equipo.}
\label{fig:escenario}
\end{figure}

\begin{figure*}[!t]
\centering
\begin{tikzpicture}[
  font=\scriptsize,
  >={Stealth[length=4pt]},
  blk/.style={draw, rounded corners=2pt, align=center, minimum height=1.15cm,
              inner sep=2pt, fill=blue!4, text width=1.6cm},
  rx/.style={blk, fill=orange!8},
  ch/.style={blk, fill=red!6, minimum height=0.85cm, text width=1.9cm},
  sum/.style={draw, circle, inner sep=0pt, minimum size=0.45cm, fill=white},
  obs/.style={draw=green!50!black, fill=green!8, rounded corners=1pt, font=\tiny,
              align=center, inner sep=1.5pt},
  tt/.style={font=\ttfamily\tiny}
]
  % Transmisor
  \node[blk] (src) {Fuente binaria\\[1pt] \texttt{\tiny Random Source}\\[-1pt]\texttt{\tiny Pack/Unpack}};
  \node[blk, right=0.35cm of src] (map) {Mapeo BPSK\\ $\{-1,+1\}$\\[1pt] \texttt{\tiny Chunks to Symbols}};
  \node[blk, right=0.35cm of map] (int) {Pulso rectangular\\ SPS $=8$\\[1pt] \texttt{\tiny Interp. FIR Filter}};
  \node[sum, right=1.0cm of int] (add) {$+$};
  % Canal
  \node[ch, above=0.45cm of add] (awgn) {AWGN de fondo\\ \texttt{\tiny noise\_amp} $=0{,}10$};
  \node[ch, below=0.45cm of add] (imp) {Tren de impulsos\\ $A_i$, $W=4$, $N_T=500$};
  % Receptor
  \node[rx, right=2.3cm of add] (mf) {Filtro acoplado\\ promedio de 8\\[1pt] \texttt{\tiny Moving Average}};
  \node[rx, right=0.35cm of mf] (smp) {Muestreo\\ 1 de cada 8\\[1pt] \texttt{\tiny Keep One in N}};
  \node[rx, right=0.35cm of smp] (dec) {Decisión\\ $\mathrm{Re}\{z_k\}\gtrless 0$\\[1pt] \texttt{\tiny Binary Slicer}};
  \node[rx, right=0.35cm of dec, text width=1.45cm] (ber) {BER\\[1pt] \texttt{\tiny XOR +}\\[-1pt]\texttt{\tiny Moving Average}};

  \draw[->] (src) -- node[above=1pt]{$R_b$} (map);
  \draw[->] (map) -- node[above=1pt]{$R_s$} (int);
  \draw[->] (int) -- node[above=1pt]{$f_s$} (add);
  \draw[->] (awgn) -- (add);
  \draw[->] (imp) -- (add);
  \draw[->] (add) -- (mf);
  \draw[->] (mf) -- (smp);
  \draw[->] (smp) -- (dec);
  \draw[->] (dec) -- (ber);

  % Bits de referencia
  \coordinate (yref) at ([yshift=-0.35cm]imp.south);
  \draw[->, dashed, gray] (src.south) -- (src.south |- yref) -- (ber.south |- yref) -- (ber.south);
  \node[above=0pt, text=black!70, font=\tiny] at ($(src.south |- yref)!0.5!(int.south |- yref)$)
       {bits transmitidos (referencia para la BER)};

  % Bloques del canal
  \begin{scope}[on background layer]
    \node[draw=red!50, dashed, rounded corners=3pt, inner sep=4pt, fit=(awgn)(imp)(add),
          label={[text=red!60!black, font=\tiny, anchor=south]north:Canal (banda base)}] {};
    \node[draw=blue!40, dashed, rounded corners=3pt, inner sep=4pt, fit=(src)(map)(int),
          label={[text=blue!60!black, font=\tiny, anchor=south]north:Transmisor}] {};
    \node[draw=orange!60, dashed, rounded corners=3pt, inner sep=4pt, fit=(mf)(smp)(dec)(ber),
          label={[text=orange!60!black, font=\tiny, anchor=south]north:Receptor}] {};
  \end{scope}

  % Puntos de observacion
  \node[obs, below=0.12cm of mf] (o2) {Ojo};
  \node[obs, below=0.12cm of smp] (o3) {Constelación\\ de decisión};
  \node[obs] at ($(add.east)!0.62!(mf.west)+(0,0.75)$) (o1) {Tiempo, PSD,\\ constelación};
  \draw[green!50!black, densely dotted] (o1.south) -- ($(add.east)!0.62!(mf.west)$);
\end{tikzpicture}
\caption{Diagrama técnico del sistema Tx--canal--Rx en banda base:
$R_s=\SI{32}{\kilo\sym\per\second}$, $R_b=\SI{32}{\kilo\bit\per\second}$,
SPS $=8$ y $f_s=\SI{256}{\kilo\samples\per\second}$. Los recuadros verdes
indican los puntos de observación. Implementación:
\texttt{S1\_BPSK\_Ruido\_Impulsivo\_Grupo3.grc} (bloques agregados en la Tabla~\ref{tab:bloques}).}
\label{fig:bloques}
\end{figure*}

\subsection{Efecto del impulso sobre la decisión BPSK}\label{subsec:efecto}

Con pulso rectangular, el filtro acoplado promedia las SPS muestras de cada
símbolo. La variable de decisión del símbolo $k$ es
\begin{equation}
  z_k=\frac{1}{\mathrm{SPS}}\sum_{n\in k} r[n] = A\,a_k + A_{\mathrm{ef}} + \tilde{w}_k,
  \label{eq:zk}
\end{equation}
donde $\tilde{w}_k$ tiene desviación $\sigma_z=\sigma_w/\sqrt{\mathrm{SPS}}$ y
$A_{\mathrm{ef}}=A_iW/\mathrm{SPS}$ es la amplitud efectiva de un impulso que cae
completo en el símbolo: el filtro acoplado lo diluye. Si una fracción
$\rho=\mathrm{SPS}/N_T$ de los símbolos contiene un impulso positivo y $p_-$ es la
fracción de ellos que vale $-1$ (1/2 con bits equiprobables),
\begin{equation}
  P_b\approx(1-\rho)\,Q\!\left(\tfrac{A}{\sigma_z}\right)
  +\rho\left[p_-Q\!\left(\tfrac{A-A_{\mathrm{ef}}}{\sigma_z}\right)
  +(1-p_-)Q\!\left(\tfrac{A+A_{\mathrm{ef}}}{\sigma_z}\right)\right].
  \label{eq:pb_imp}
\end{equation}
Con $\sigma_z\ll A$, la BER es nula si $A_i<A\,\mathrm{SPS}/W$, vale
$\rho p_-/2$ en el umbral y alcanza el piso $\rho p_-$ por encima de él:
\begin{equation}
  P_b\approx 0 \;\;(A_i<A\,\mathrm{SPS}/W),\qquad
  P_b\approx\rho\,p_- \;\;(A_i>A\,\mathrm{SPS}/W).
  \label{eq:pb_limite}
\end{equation}
El comportamiento es de umbral y el piso no depende de la energía de la señal,
como se ha reportado para sistemas de portadora única~\cite{ghosh1996}. Como
comparación, un receptor que tomara una sola muestra por símbolo, sin filtro
acoplado, tendría el umbral en $A_i=A$, pero solo fallaría cuando el impulso
coincide con el instante de muestreo, con $\rho\approx W/N_T$. El filtro acoplado
eleva entonces el umbral de amplitud a cambio de que todo impulso influya en
alguna decisión. Si el impulso tuviera componente en cuadratura, dispersaría la
constelación en vertical sin alterar la decisión BPSK, que depende solo de la
parte real: la dispersión visible no equivale necesariamente a errores.

%============================================================================
\section{Diseño del sistema y parámetros}\label{sec:diseno}
%============================================================================

\subsection{Escenario y parámetros}\label{subsec:escenario}

Un nodo sobre una máquina transmite variables de baja velocidad a la estación
de supervisión, en una nave con motores, contactores y variadores cuyas
conmutaciones generan perturbaciones breves de gran amplitud
(Fig.~\ref{fig:escenario}).

Se verificó que el paquete de referencia del curso fuera adecuado: una
vibración muestreada a \SI{1}{\kilo\samples\per\second} con 16 bits requiere
\SI{16}{\kilo\bit\per\second}, la mitad de la tasa disponible, y la temperatura
o el estado requieren mucho menos. Las relaciones de diseño son
\begin{align}
  k &= \log_2 2 = 1, \qquad R_b = R_s\,k = \SI{32}{\kilo\bit\per\second}, \label{eq:rb}\\
  f_s &= R_s\cdot\mathrm{SPS} = \num{32000}\times 8 = \SI{256}{\kilo\samples\per\second}, \label{eq:fs}\\
  B_{\mathrm{nulos}} &\approx 2R_s=\SI{64}{\kilo\hertz}, \label{eq:bw}
\end{align}
coherentes con el flujograma base. La ventana de la PSD
($\pm f_s/2=\pm\SI{128}{\kilo\hertz}$) contiene el lóbulo principal y tres
lóbulos laterales por lado. La Tabla~\ref{tab:parametros} resume el diseño.

\begin{table}[!t]
\centering
\caption{Parámetros de diseño del sistema}
\label{tab:parametros}
\footnotesize
\begin{tabular}{@{}llr@{}}
\toprule
\textbf{Parámetro} & \textbf{Símbolo} & \textbf{Valor} \\
\midrule
Modulación              & --                     & BPSK ($M=2$) \\
Bits por símbolo        & $k$                    & 1 \\
Tasa de símbolos        & $R_s$                  & \SI{32}{\kilo\sym\per\second} \\
Tasa de bits            & $R_b$                  & \SI{32}{\kilo\bit\per\second} \\
Periodo de símbolo      & $T_s=T_b$              & \SI{31,25}{\micro\second} \\
Muestras por símbolo    & SPS                    & 8 \\
Frecuencia de muestreo  & $f_s$                  & \SI{256}{\kilo\samples\per\second} \\
Conformación de pulso   & $p(t)$                 & Rectangular \\
Ancho de banda (nulos)  & $B_{\mathrm{nulos}}$   & \SI{64}{\kilo\hertz} \\
Ancho de banda RRC (ref.)& $B_{\mathrm{RRC}}$    & \SI{43,2}{\kilo\hertz} \\
Eficiencia espectral    & $R_b/B_{\mathrm{nulos}}$ & \SI{0,5}{\bit\per\second\per\hertz} \\
Perturbación            & $i[n]$                 & Impulsiva periódica \\
Variable experimental   & $A_i$                  & Amplitud del impulso \\
\bottomrule
\end{tabular}
\end{table}

\subsection{Arquitectura del sistema}\label{subsec:arquitectura}

En la Fig.~\ref{fig:bloques}, la fuente binaria alimenta el mapeo BPSK y el
interpolador de pulso rectangular; el canal suma un AWGN de fondo fijo y el
tren de~\eqref{eq:tren}; el receptor aplica el filtro acoplado, toma una
muestra por símbolo, decide por signo y compara con los bits transmitidos
para obtener la BER.

%============================================================================
\section{Metodología de simulación}\label{sec:metodologia}
%============================================================================

\begin{figure}[!t]
\centering
\captura[\columnwidth]{grc_base.png}
\caption{Flujograma base del curso con fuente de AWGN.}
\label{fig:grc_base}
\end{figure}

\begin{figure}[!t]
\centering
\captura[\columnwidth]{grc_modificado.png}
\caption{Flujograma modificado con generador de impulsos y receptor BPSK.}
\label{fig:grc_mod}
\end{figure}

\begin{figure*}[!t]
\centering
\setlength{\tabcolsep}{1pt}
\begin{tabular}{cccc}
\captura[0.24\textwidth]{tiempo_A.png} & \captura[0.24\textwidth]{tiempo_B.png} &
\captura[0.24\textwidth]{tiempo_C.png} & \captura[0.24\textwidth]{tiempo_D.png} \\
\scriptsize (a) A: $A_i=0$ & \scriptsize (b) B: $A_i=1$ &
\scriptsize (c) C: $A_i=2$ & \scriptsize (d) D: $A_i=4$ \\[2pt]
\captura[0.24\textwidth]{psd_A.png} & \captura[0.24\textwidth]{psd_B.png} &
\captura[0.24\textwidth]{psd_C.png} & \captura[0.24\textwidth]{psd_D.png} \\
\scriptsize (e) A: $A_i=0$ & \scriptsize (f) B: $A_i=1$ &
\scriptsize (g) C: $A_i=2$ & \scriptsize (h) D: $A_i=4$ \\
\end{tabular}
\caption{Señal en el canal (componentes I y Q; amplitud frente a tiempo en ms,
ventana de \SI{8}{\milli\second}; la escala vertical se ajustó en cada prueba)
en (a)--(d), y PSD (dB frente a frecuencia en kHz, $\pm\SI{128}{\kilo\hertz}$)
en (e)--(h), para las cuatro pruebas.}
\label{fig:tiempo_psd}
\end{figure*}

\begin{figure*}[!t]
\centering
\setlength{\tabcolsep}{1pt}
\begin{tabular}{cccc}
\captura[0.24\textwidth]{ojo_A.png} & \captura[0.24\textwidth]{ojo_B.png} &
\captura[0.24\textwidth]{ojo_C.png} & \captura[0.24\textwidth]{ojo_D.png} \\
\scriptsize (a) A: $A_i=0$ & \scriptsize (b) B: $A_i=1$ &
\scriptsize (c) C: $A_i=2$ & \scriptsize (d) D: $A_i=4$ \\[2pt]
\captura[0.24\textwidth]{const_A.png} & \captura[0.24\textwidth]{const_B.png} &
\captura[0.24\textwidth]{const_C.png} & \captura[0.24\textwidth]{const_D.png} \\
\scriptsize (e) A: $A_i=0$ & \scriptsize (f) B: $A_i=1$ &
\scriptsize (g) C: $A_i=2$ & \scriptsize (h) D: $A_i=4$ \\
\end{tabular}
\caption{Diagrama de ojo a la salida del filtro acoplado (componente I; amplitud
frente a tiempo en ms) en (a)--(d), y constelación en el instante de decisión
(I frente a Q) en (e)--(h), para las cuatro pruebas.}
\label{fig:ojo_const}

\vspace{6pt}
\setlength{\tabcolsep}{1pt}
\begin{tabular}{cccc}
\captura[0.24\textwidth]{const_sobre_A.png} & \captura[0.24\textwidth]{const_sobre_B.png} &
\captura[0.24\textwidth]{const_sobre_C.png} & \captura[0.24\textwidth]{const_sobre_D.png} \\
\scriptsize (a) A: $A_i=0$ & \scriptsize (b) B: $A_i=1$ &
\scriptsize (c) C: $A_i=2$ & \scriptsize (d) D: $A_i=4$ \\
\end{tabular}
\caption{Constelación sobremuestreada a la salida del canal (todas las
muestras, I frente a Q) para las cuatro pruebas.}
\label{fig:const_sobre}
\end{figure*}

\subsection{Flujograma base y modificaciones}\label{subsec:grc_base}

Se partió de \texttt{Modulaciones\_basadas\_en\_constelaciones\_AWGN.grc}, que
conserva el transmisor del archivo base y suma ruido gaussiano
(Fig.~\ref{fig:grc_base}). Su fuente aleatoria repite una secuencia de 1000
bits, el bloque \texttt{Chunks to Symbols} usa la tabla de constelación y el
interpolador con \texttt{h=[1]*Sps} forma el pulso rectangular. El original se
conservó y se trabajó sobre la copia
\texttt{S1\_BPSK\_Ruido\_Impulsivo\_Grupo3.grc}. Como la tabla por defecto era
la de OOK, $(0,1)$, se cambió por la de BPSK, $(-1,+1)$, y se verificó que
$M=2$ y $R_b=\SI{32}{\kilo\bit\per\second}$.

La Tabla~\ref{tab:bloques} y la Fig.~\ref{fig:grc_mod} muestran los bloques
agregados. El tren de impulsos se generó con un \texttt{Vector Source} de
$N_T$ muestras, $W$ de ellas en uno, multiplicado por \texttt{A\_imp} y sumado
en una tercera entrada del \texttt{Add}. El receptor usa un promedio móvil de
SPS muestras como filtro acoplado, un \texttt{Keep One in N} que conserva la
última muestra de cada símbolo (máximo del integrador) y un
\texttt{Binary Slicer}, que entrega 1 si la parte real es positiva, de acuerdo
con la asignación $0\to-1$ y $1\to+1$. La BER se obtiene con una XOR entre bits
transmitidos y decididos, cuya salida vale 1 en cada error, promediada sobre
$10^5$ bits (unos \SI{3,1}{\second}). Como ni el interpolador ni el promedio
móvil introducen retardo de símbolo, ambas secuencias quedan alineadas; la
alineación se validó con BER nula en la referencia.

\begin{table}[!t]
\centering
\caption{Bloques agregados al flujograma}
\label{tab:bloques}
\scriptsize
\setlength{\tabcolsep}{3pt}
\begin{tabular}{@{}p{0.30\columnwidth}p{0.30\columnwidth}p{0.34\columnwidth}@{}}
\toprule
\textbf{Bloque} & \textbf{Parámetros} & \textbf{Función} \\
\midrule
\texttt{Vector Source} (complex) & \texttt{[1]*W\_imp + [0]*(N\_T-W\_imp)}, repetir & Forma del tren de impulsos \\
\texttt{Multiply Const} & \texttt{A\_imp} & Amplitud del impulso (variable) \\
\texttt{Add} (3 entradas) & -- & Señal + AWGN + impulsos \\
\texttt{Moving Average} (complex) & longitud 8, escala $1/8$ & Filtro acoplado \\
\texttt{QT GUI Eye Sink} & 8 muestras/símbolo & Diagrama de ojo \\
\texttt{Keep One in N} & $N=8$ & Una muestra por símbolo \\
\texttt{QT GUI Constellation Sink} & -- & Constelación de decisión \\
\texttt{Complex to Real} + \texttt{Binary Slicer} & umbral 0 & Decisión de bit \\
\texttt{XOR} + \texttt{Char to Float} & -- & Detección de errores \\
\texttt{Moving Average} (float) & longitud $10^5$ & Estimación de BER \\
\texttt{QT GUI Number Sink} & -- & Visualización de BER \\
\bottomrule
\end{tabular}
\end{table}

\subsection{Plan de pruebas}\label{subsec:plan}

Se fijaron $W=4$ muestras (\SI{15,6}{\micro\second}, medio símbolo) y
$N_T=500$ (\SI{1,953}{\milli\second}). La repetición es más rápida que la de una
conmutación real para observar cuatro impulsos por captura y reunir errores en
pocos segundos. Como 500 no es múltiplo de 8, los impulsos caen alternadamente
en la primera y la segunda mitad de un símbolo, pero siempre completos dentro
de él, de modo que
$\rho=\mathrm{SPS}/N_T=0{,}016$ y $A_{\mathrm{ef}}=A_i/2$. El AWGN de fondo
(\texttt{noise\_amp} $=0{,}10$, $\sigma_z=0{,}025$) no produce errores por sí
solo. Solo cambió $A_i$, con niveles alrededor del umbral teórico
$A\,\mathrm{SPS}/W=2$ (Tabla~\ref{tab:pruebas}, con
$\mathrm{SIR}=A^2/(A_i^2\delta)$ y la BER esperada para bits equiprobables). En cada prueba se esperaron
al menos \SI{5}{\second} antes de registrar las gráficas y la BER.

\begin{table}[!t]
\centering
\caption{Plan de pruebas: condición de referencia y tres niveles de perturbación}
\label{tab:pruebas}
\footnotesize
\setlength{\tabcolsep}{4pt}
\begin{tabular}{@{}clcccc@{}}
\toprule
\textbf{Prueba} & \textbf{Severidad} & \textbf{$A_i$} & \textbf{$A_{\mathrm{ef}}$} &
\textbf{SIR [dB]} & \textbf{BER esperada} \\
\midrule
A & Referencia & 0   & 0   & $\infty$ & 0 \\
B & Suave      & 1,0 & 0,5 & 21,0 & $\approx 0$ \\
C & Media      & 2,0 & 1,0 & 14,9 & $\approx 4\times10^{-3}$ \\
D & Fuerte     & 4,0 & 2,0 & 8,9  & $\approx 8\times10^{-3}$ \\
\midrule
\multicolumn{6}{@{}p{0.95\columnwidth}@{}}{Constantes: $W=4$, $N_T=500$,
\texttt{noise\_amp} $=0{,}10$, $N_{\mathrm{ber}}=10^5$; $R_s$, SPS y $f_s$ según la
Tabla~\ref{tab:parametros}.} \\
\bottomrule
\end{tabular}
\end{table}

\subsection{Hipótesis}\label{subsec:hipotesis}

Antes de simular se esperaba que mayores amplitudes produjeran perturbaciones
temporales más fuertes y mayor dispersión. Con el modelo anterior: impulsos
visibles desde B sin errores; en C, BER cercana a $\rho/4=4\times10^{-3}$; en D,
un piso de $\rho/2=8\times10^{-3}$ independiente de la potencia de la señal; y
en la PSD, un nivel adicional de banda ancha creciente con $A_i^2$.
